\documentclass[../../main.tex]{subfiles}
\begin{document}

{\bf [解]}

題意で与えられた$x_i$の条件
\begin{align}
x_i>0, \qquad \sum_{i=1}^nx_i=k \quad\cdots\text{①} \label{eq:1}
\end{align}
の元で考える．簡単のため
\begin{align}
 \begin{cases}
  A_n&=\displaystyle\sum_{i=1}^nx_i\log x_i \\
  B_n&=k\log\dfrac{k}{n}
 \end{cases}
\end{align}
とおく．「$A_n\ge B_n\ \cdots$が任意の$n\in\mathbb{N}$で成立すること」を帰納法で示す．
$n=1$の時には$A_1=B_1=k\log k$だから成立は明らかなので，以下$n=p\, (p\in\mathbb{N})$での成立を仮定して$n=p+1$でも成立することを示す．

$p+1$個の正数$x_i\ (i=1,2,\cdots,p+1)$に対し，$\displaystyle\sum_{i=1}^{p+1}x_i=k$とし，$\displaystyle\sum_{i=1}^px_i=\ell\ (0<\ell<k)$とおく．まず$\ell$を固定して考えると仮定より
\begin{align}
\begin{split}
A_{p+1}
 &=A_p+x_{p+1}\log x_{p+1} \\
 &\ge\ell\log\dfrac{\ell}{p}+(k-\ell)\log(k-\ell)\equiv f(\ell) \quad(\because\text{仮定}) \quad\cdots\text{②} \label{eq:2}
\end{split}
\end{align}
と下から評価できる．次に$\ell$を動かして$f(\ell)$の最小値を考えよう．
$f$の一階微分を計算すると
\begin{align}
f'(\ell)=1+\log\ell-\log p-\log(k-\ell)-1=\log\dfrac{\ell}{k-\ell}-\log p
\end{align}
であって，$\log x$が$x$の単調増加関数であるから増減表は以下のようになる．
\begin{table}[htb]
  \centering
\begin{tabular}{c|ccccc}
$\ell$ & $0$ & & $\dfrac{pk}{p+1}$ & & $k$ \\
\hline
$f'$ & & $-$ & $0$ & $+$ & \\
\hline
$f$ & & $\searrow$ & & $\nearrow$ &
\end{tabular}
\end{table}

従って$f$は下から評価できて
\begin{align}
\begin{split}
f(\ell)
 &\ge f\left(\dfrac{pk}{p+1}\right) \\
 & = \dfrac{pk}{p+1}\log\dfrac{k}{p+1}+\dfrac{k}{p+1}\log\dfrac{k}{p+1} \\
 & = k\log\dfrac{k}{p+1}=B_{p+1} \quad\cdots\text{③} \label{eq:3}
\end{split}
\end{align}
\cref{eq:3}を\cref{eq:2}に代入して
\begin{align}
 A_{p+1} \ge B_{p+1}
\end{align}
である．よって$n=p+1$でも主張は成立する．

以上から数学的帰納法により$A_n \ge B_n$であり，題意は示された．


\newpage

\end{document}
